% Prime-Force Correspondence
\documentclass[aps,prl,twocolumn,superscriptaddress,nofootinbib]{revtex4-2}

\usepackage{amsmath,amssymb,amsfonts}
\usepackage{hyperref}
\usepackage{booktabs}

\newcommand{\Spec}{\mathrm{Spec}}
\newcommand{\Gal}{\mathrm{Gal}}
\newcommand{\reg}{\mathrm{reg}}
\newcommand{\SU}{\mathrm{SU}}
\newcommand{\U}{\mathrm{U}}

\begin{document}

\title{Prime-Force Correspondence:\\
Gauge Groups from Ramified Coverings of $\Spec(\mathbb{Z})$}

\author{Wright Brothers Research Program}
\affiliation{Independent Research, 2026}

\date{\today}

\begin{abstract}
We observe that the three gauge group factors of the Standard Model
admit a uniform description in terms of number field extensions of
$\mathbb{Q}$ ramified at a single prime:
$\U(1)\leftrightarrow\mathbb{Q}$ (trivial),
$\SU(2)\leftrightarrow\mathbb{Q}(\sqrt{2})$ (ramified at $p{=}2$),
$\SU(3)\leftrightarrow\mathbb{Q}(\zeta_9)^+$ (ramified at $p{=}3$).
The pattern is $\SU(N)\leftrightarrow$ degree-$N$ covering with
Galois group $\mathbb{Z}/N$, ramified at $p{=}N$.  Two independent
numerical coincidences support this correspondence:
$\cos\theta_W=\reg(\mathbb{Q}(\sqrt{2}))=\log(1{+}\sqrt{2})$
agrees with $m_W/m_Z$ to $0.001\%$, and
$10\,h\!\cdot\!R(\mathbb{Q}(\zeta_9)^+)=8.493$ agrees with
$\alpha_s(M_Z)^{-1}=8.482$ to $0.13\%$.  We propose that gauge
symmetry breaking is encoded in the ramification structure of
arithmetic coverings, with the Frobenius element playing the role
of gauge holonomy.  The correspondence makes a falsifiable prediction
for $m_W$ at sub-MeV precision.
\end{abstract}

\maketitle

%=====================================================================
\section{Observation}
%=====================================================================

The Standard Model gauge group $\SU(3)_c\times\SU(2)_L\times\U(1)_Y$
has resisted derivation from a deeper principle.  We report a
numerical observation that may point toward such a principle.

Consider the following number fields:
\begin{itemize}
\item $\mathbb{Q}$ itself (degree~1 over $\mathbb{Q}$).
\item $K_2=\mathbb{Q}(\sqrt{2})$, the simplest real quadratic field.
  Discriminant $\Delta=8=2^3$; ramified at $p=2$ only; class number
  $h=1$; regulator $R=\log(1{+}\sqrt{2})=0.88137\ldots$
\item $K_3=\mathbb{Q}(\zeta_9)^+=\mathbb{Q}(\cos 2\pi/9)$, the
  maximal real subfield of the 9th cyclotomic field.  Degree~3;
  discriminant $\Delta=81=3^4$; ramified at $p=3$ only; Galois group
  $\mathbb{Z}/3$; $hR=0.84931\ldots$
\end{itemize}

Each field is the \emph{minimal} extension of $\mathbb{Q}$ with
cyclic Galois group $\mathbb{Z}/N$ that is ramified at $p=N$
\emph{only}, for $N=1,2,3$ respectively.

We observe:

\begin{equation}\label{eq:cosW}
  \cos\theta_W = \frac{m_W}{m_Z} \stackrel{?}{=} \reg(K_2)
  = \log(1{+}\sqrt{2}).
\end{equation}

\noindent Numerically: $\log(1{+}\sqrt{2})=0.881374$ vs.\
$m_W/m_Z=0.881361\pm 0.000147$.  Agreement: $0.09\sigma$.

\begin{equation}\label{eq:alphas}
  \frac{1}{\alpha_s(M_Z)} \stackrel{?}{=} 10\;hR(K_3)
  = 10\times 0.84931.
\end{equation}

\noindent Numerically: $8.493$ vs.\ $1/\alpha_s=8.482\pm 0.008$.
Agreement: $0.13\%$.

%=====================================================================
\section{The Correspondence}
%=====================================================================

\begin{center}
\begin{tabular}{cccccc}
\toprule
Gauge & Field & Degree & $\Gal$ & Ram.\ at & Invariant \\
\midrule
$\U(1)$ & $\mathbb{Q}$ & 1 & $\{1\}$ & --- & $\alpha$ \\
$\SU(2)$ & $K_2$ & 2 & $\mathbb{Z}/2$ & $p{=}2$ & $\cos\theta_W$ \\
$\SU(3)$ & $K_3$ & 3 & $\mathbb{Z}/3$ & $p{=}3$ & $\alpha_s^{-1}$ \\
\bottomrule
\end{tabular}
\end{center}

The pattern: $\SU(N)\leftrightarrow$ degree-$N$ cyclic extension of
$\mathbb{Q}$, ramified at $p=N$ only, with Galois group
$\mathbb{Z}/N$.

\medskip
\noindent\textbf{Why $p=2$ and weak force?}
The weak interaction is the \emph{only} force that violates parity
($P$)~\cite{Wu1957}.  Parity is a $\mathbb{Z}/2$ symmetry, and
$p=2$ is the unique prime governing the even/odd distinction.
$K_2=\mathbb{Q}(\sqrt{2})$ is the unique class-number-one real
quadratic field ramified at $p=2$ only.  The chain
$P\text{-violation}\leftrightarrow\mathbb{Z}/2\leftrightarrow
p{=}2\leftrightarrow K_2$ provides a structural (if not yet
rigorous) link between the weak force and the arithmetic of
$\Spec(\mathbb{Z})$.

\medskip
\noindent\textbf{Why degree~3 for $\SU(3)$?}
A field ramified at $p=3$ only \emph{cannot} be quadratic: for
$d=3$, the discriminant is $\Delta=12=4\times 3$, which is ramified
at both $p=2$ and $p=3$.  This is because $3\equiv 3\pmod{4}$;
for any $p\equiv 3\pmod{4}$, the quadratic field
$\mathbb{Q}(\sqrt{p})$ has $\Delta=4p$, inevitably introducing
$p=2$.  A ``pure'' covering ramified at $p=3$ only requires
degree~$\geq 3$.  The minimal such is the cyclic cubic
$K_3=\mathbb{Q}(\zeta_9)^+$, with $\Delta=3^4$.  Thus
\emph{number theory itself forces} $\SU(3)$ to live in a
higher-degree covering than $\SU(2)$.

%=====================================================================
\section{The Class Number Formula Connection}
%=====================================================================

For $K_2$, the Dirichlet class number formula gives:
\begin{equation}\label{eq:classK2}
  L(\chi_8, 1) = \frac{hR}{\sqrt{\Delta}}
  = \frac{\log(1{+}\sqrt{2})}{\sqrt{8}}.
\end{equation}
Since $\cos\theta_W=hR=\log(1{+}\sqrt{2})$ and $h=1$:
\begin{equation}
  \cos\theta_W = L(\chi_8, 1)\,\sqrt{\Delta(K_2)}.
\end{equation}

For $K_3$, the class number formula gives $hR$ in terms of
$|L(\chi_3,1)|^2$ where $\chi_3$ is a character of order~3 mod~9:
\begin{equation}
  hR(K_3) = \tfrac{9}{4}\,|L(\chi_3,1)|^2.
\end{equation}

In both cases, the physical parameter (mixing angle or coupling)
equals an arithmetic invariant that appears in the class number
formula---one of the deepest results in algebraic number
theory~\cite{Dirichlet1837}.

The Dedekind zeta function of $K$ factorizes as
$\zeta_K(s)=\zeta(s)\cdot L(\chi,s)\cdots$, separating into the
``base'' zeta (governing $\alpha$ and $G$) and the $L$-factors
(governing mixing angles and non-abelian couplings).  At $s=1$,
$\zeta(s)$ has a pole (gravity), while each $L$-factor is finite
(mixing parameters).

%=====================================================================
\section{Frobenius as Holonomy}
%=====================================================================

In the covering $K_2/\mathbb{Q}$, each unramified prime $p\neq 2$
carries a Frobenius element $\mathrm{Frob}_p\in\Gal(K_2/\mathbb{Q})
=\mathbb{Z}/2$:

\begin{center}
\begin{tabular}{ccc}
\toprule
Prime type & $\mathrm{Frob}_p$ & Physical analog \\
\midrule
Split ($p\equiv\pm 1\!\pmod{8}$) & $\mathrm{id}$ & trivial holonomy \\
Inert ($p\equiv\pm 3\!\pmod{8}$) & $\sigma$ & nontrivial holonomy \\
Ramified ($p=2$) & undefined & source of the force \\
\bottomrule
\end{tabular}
\end{center}

In gauge theory, the holonomy of a connection around a loop
classifies the local gauge transformation.  We propose that the
Frobenius element at each prime $p$ plays the role of holonomy of
an arithmetic gauge connection, with the ramification locus ($p=2$
for the weak force) as the ``source'' of the gauge field.

%=====================================================================
\section{Predictions}
%=====================================================================

\noindent\textbf{1.}
$m_W = m_Z\times\log(1{+}\sqrt{2}) = 80.370\;\mathrm{GeV}$.
Current experimental value: $80.369\pm 0.013\;\mathrm{GeV}$ (PDG
2024), $0.09\sigma$ from prediction.  Future measurements at
sub-MeV precision (FCC-ee, CEPC) will provide a decisive test.

\noindent\textbf{2.}
$\sin^2\theta_W = 1-\log^2(1{+}\sqrt{2}) = 0.22318$, to be
compared with the on-shell value $0.22320\pm 0.00004$.

\noindent\textbf{3.}
The strong coupling is related to $K_3$ via $\alpha_s^{-1}(M_Z)
= c\cdot hR(K_3)$ with $c=10$ (the origin of this factor is
currently unexplained).

%=====================================================================
\section{What This Is and What It Is Not}
%=====================================================================

This paper reports a \emph{numerical observation} and a
\emph{structural correspondence}, not a derivation.  We do not
derive the Standard Model gauge group from number theory; we
observe that the group structure matches the ramification structure
of specific number field extensions in a way that produces correct
physical parameters.

The key weakness is the absence of a dynamical mechanism: we do not
explain \emph{why} gauge couplings should equal arithmetic
regulators.  The factor of~10 in the $\alpha_s$ relation is
unexplained.  The correspondence does not extend to CKM or PMNS
mixing matrices.

The key strength is falsifiability:
Eq.~\eqref{eq:cosW} predicts $m_W$ to $0.1\;\mathrm{MeV}$, well
within reach of next-generation colliders.  If $m_W$ deviates from
$m_Z\log(1{+}\sqrt{2})$ by more than a few MeV, the correspondence
is ruled out.

The correspondence $\SU(N)\leftrightarrow$ ramification at $p=N$
is forced by the structure of algebraic number theory (the
impossibility of a quadratic field ramified at $p=3$ only).
Whether nature respects this arithmetic constraint is an empirical
question.

\begin{thebibliography}{99}
\bibitem{Wu1957} C.~S.~Wu \textit{et al.}, Phys.\ Rev.\ \textbf{105}, 1413 (1957).
\bibitem{Dirichlet1837} P.~G.~L.~Dirichlet, Abh.\ K\"onigl.\ Preuss.\ Akad.\ Wiss.\ (1837).
\bibitem{Neukirch1999} J.~Neukirch, \textit{Algebraic Number Theory} (Springer, 1999).
\bibitem{Connes1996} A.~Connes, Commun.\ Math.\ Phys.\ \textbf{182}, 155 (1996).
\end{thebibliography}

\end{document}
